\documentclass[10pt,a4paper]{amsart}
\usepackage[margin=19mm]{geometry}
\usepackage{amsmath,amssymb,mathtools}
\usepackage[scr=rsfs,cal=euler]{mathalfa}
\usepackage{xfrac}
\usepackage[unicode,hidelinks,bookmarks=false]{hyperref}
\newcommand{\ie}{i.e.\ }
\newcommand{\cf}{cf.\ }
\newcommand{\resp}{respectively\ }
\newcommand{\Subsection}{\subsection}
\providecommand{\nobreakdash}{\nobreak\hbox{-}}
\setlength{\emergencystretch}{2em}
\multlinegap0pt
\usepackage{bm}
\usepackage{bbm}
\usepackage{dsfont}
\usepackage{xspace}
\usepackage{cancel}
\usepackage{mathtools}
\usepackage{amsthm}
\makeatletter
\@ifundefined{theorem}{\newtheorem{theorem}{Theorem}}{}
\@ifundefined{lemma}{\newtheorem{lemma}[theorem]{Lemma}}{}
\makeatother
\title[Characterisation of Convergence, Boundedness and Unboundedne]{Characterisation of Convergence, Boundedness and Unboundedness in Solutions of Second Order Linear Differential Equations}
\author{John A. D. Appleby, Subham Pal}
\date{Source: arXiv:2603.25545v1 (CC BY 4.0)}
\begin{document}
\maketitle

\noindent\emph{Source and licence.} Characterisation of Convergence, Boundedness and Unboundedness in Solutions of Second Order Linear Differential Equations. Version arXiv:2603.25545v1 (CC BY 4.0), distributed under the Creative Commons Attribution 4.0 International licence, as recorded in the arXiv licence metadata for this exact version and shown on the abstract page https://arxiv.org/abs/2603.25545v1. Full rights notice: licenses/arxiv-2603.25545-CC-BY-4.0.txt.\par\smallskip

\noindent\textit{Editorial scope.} Counted unit: the Lemma whose source label is \texttt{lemma:delta\_integral} in the article (source0.tex lines 376--386, source label \texttt{lemma:delta\_integral}), delivered together with the article's own complete original proof of it (source0.tex lines 411--489, one \texttt{proof} environment of 79 lines). No mathematics is written, repaired or re-derived: statement and proof are the article's own text line for line. The proof is detached from the statement in the source: the article states the result at lines 376--386 and prints its proof at lines 411--489, and the proof environment's own header names the statement, so the association is the article's own. Required context retained uncounted: none. Every definition, display and label of those lines is retained unchanged. Omitted from this package and not redistributed: 1--375, 387--410, 490--855. Nothing inside a retained excerpt is omitted or altered: every retained body is delivered exactly as the article's lines are written, so no omission receipt of any kind accompanies this package. Citations: the retained excerpts carry no citation command at all, so no bibliography entry of the article is transcribed and no reference is redistributed. Reference adaptation: none was needed and none was made; every reference inside the delivered unit resolves inside it, so no unresolved reference or citation is produced. Printed slips kept literal: no printed word, symbol, hypothesis or number of the counted statement or of its proof was altered. Layout and class: the article's own class is \texttt{amsart} with packages including amsmath, amssymb, amsthm, geometry; the wrapper is the lane's pinned \texttt{amsart} editorial wrapper at 10pt on A4, which restates the theorem-like environments the retained text needs and adds the article's own macro definitions that the retained text needs (none), with extra wrapper packages (bm, bbm, dsfont, xspace, cancel, mathtools). The delivered numbering is the unit's own. Assets: no retained span contains a figure environment, a graphics-inclusion command, a PDF-page inclusion, a table or a TikZ picture; no image file is copied into this package, the package declares no asset dependency and no image file is redistributed with it. Licence: version arXiv:2603.25545v1 of this article is distributed under the Creative Commons Attribution 4.0 International licence, as recorded in the arXiv licence metadata for this exact version and shown on the abstract page https://arxiv.org/abs/2603.25545v1. A full rights notice is delivered at licenses/arxiv-2603.25545-CC-BY-4.0.txt. Characters: every retained excerpt is pure ASCII, so the delivered engine, XeLaTeX with its default Latin Modern text font, reports no missing character.
	\begin{lemma} \label{lemma:delta_integral}
		Suppose $g \in L_{loc}^1[0, \infty)$. For each $\theta \in [0, 1]$, define the moving average integral:
		\[
		g_\theta(t) = \int_{(t-\theta)^+}^t g(s) \, ds, \quad t \ge 0.
		\]
		Define the difference function $(\delta g)(t) = g(t) - g_1(t)$.
		Then for all $t \ge 0$:
		\begin{equation}
			\int_0^t (\delta g)(s) \, ds = \int_0^1 g_\theta(t) \, d\theta.
		\end{equation}
	\end{lemma}

	\begin{proof}
		We start with $I = k * f$ and decompose $f = f_1 + \delta f$: write
		\[ I = I_1 + I_2, \quad \text{where } I_1 = k * f_1, \quad I_2 = k * \delta f. \]
		\textbf{Step 1: Analysis of $I_2$}
		We apply integration by parts to $I_2 = \int_0^t k(t-s) (\delta f)(s) \, ds$. Using Lemma \ref{lemma:delta_integral}, we have that $b(s):=\int_0^1 f_\theta(s)\,d\theta$ implies $b'(s)=(\delta f)(s)$. Therefore 
		\[
		\int_0^t k(t-s)(\delta f)(s)\,ds
		= k(0)b(t)
		-k(t)b(0)+\int_0^t k'(t-s)b(s)\,ds
		=(k'\ast b)(t).
		\]
		Now, using the notation of Lemma~\ref{lemma:delta_integral}, we write $b = b_1 + \delta b$, so that 
		\begin{equation} \label{eq.I2}
			I_2(t)=\int_0^t k'(t-s)b_1(s)\,ds+\int_0^t k'(t-s)(\delta b)(s)\,ds.
		\end{equation}
		We now analyse these two integrals. First, we find $b_1$. Applying the definitions of $b_1$, $f_\theta$, and $F_{(1,\theta)}$ in order, we get  
		\begin{equation} \label{eq.b1}
			b_1(t) = \int_{(t-1)^+}^t b(s) \, ds = \int_0^1 \int_{(t-1)^+}^t f_\theta(s) \, ds \, d\theta
			=\int_0^1 \int_{(t-1)^+}^t
			\int_{(s-\theta)^+}^s f(u) \, du
			\,ds\,d\theta
			= \int_0^1 F_{(1,\theta)}(t)\,d\theta.
		\end{equation}
		Next, we deal with $(k'\ast (\delta b))(t)$. Using the definition of $b_\phi(t):=\int_{(t-\phi)^+}^t b(s)\,ds$ and applying Lemma \ref{lemma:delta_integral}, we have $\int_0^s (\delta b)(u)\,du = \int_0^1 b_\phi(s)\,d\phi$. Thus, if $v(s)=\int_0^1 b_\phi(s)\,d\phi$, then $v'(s)=(\delta b)(s)$. Now, integrating by parts, we get  
		\begin{equation} \label{eq.kprdelb}
			\int_0^t k'(t-s)(\delta b)(s)\,ds
			=k'(0)v(t)-k'(t)v(0)+\int_0^t k''(t-s)v(s)\,ds.
		\end{equation}
		Now, to compute $v$, we need $b_\phi$. Using in turn the definition of $b_\phi$, then the definition of $b$, Fubini's theorem, and finally the definitions of $f_\theta$ and $F$, we arrive at 
		\[ b_\phi(s) = \int_{(s-\phi)^+}^s b(u) \, du = \int_{(s-\phi)^+}^s \int_0^1 f_\theta(s)\,d\theta\,du = \int_0^1 \int_{(s-\phi)^+}^s f_\theta(u) \, du \, d\theta = \int_0^1 F_{(\phi, \theta)}(s) \, d\theta. \]
		Therefore
		\begin{equation} \label{eq.v}
			v(t)=\int_0^1 b_\phi(t)\,d\phi = \int_0^1 \int_0^1 F_{(\phi,\theta)}(t)\,d\theta\, d\phi.
		\end{equation}
		Since $F_{(\phi,\theta)}(0)=0$, we have $v(0)=0$. 
		Therefore from \eqref{eq.kprdelb}, and using the fact that $k'(0)=1$, we have  
		\begin{equation*} 
			\int_0^t k'(t-s)(\delta b)(s)\,ds
			=v(t)+\int_0^t k''(t-s)v(s)\,ds.
		\end{equation*}
		Combining this with \eqref{eq.v}, 
		\eqref{eq.b1} and \eqref{eq.I2}, we get that 
		\begin{equation} \label{eq.I2final}
			I_2(t)=\left(k'\ast \int_0^1 F_{(1,\theta)}(\cdot)\,d\theta\right)(t)+ \int_0^1 \int_0^1 F_{(\phi,\theta)}(t)\,d\theta\, d\phi
			+\left(k''\ast \int_0^1 \int_0^1 F_{(\phi,\theta)}(\cdot)\,d\theta\, d\phi\right)(t).
		\end{equation}
		
		\textbf{Step 2: Analysis of $I_1$}
		It remains to consider $I_1(t) = (k * f_1)(t)$. Let $a = f_1$. Using the notation of Lemma~\ref{lemma:delta_integral}, we write $a = a_1 + \delta a$, so that 
		\[
		I_1(t)=\int_0^t k(t-s)a_1(s)\,ds + \int_0^t k(t-s)(\delta a)(s)\,ds. 
		\]
		For the first term on the right hand side, we need an expression for $a_1$. From the definitions of $a_1$, $a$, $f_1$ and $F$, we find that  
		\[
		a_1(t) = \int_{(t-1)^+}^t a(s) \, ds = \int_{(t-1)^+}^t \int_{(s-1)^+}^s f(u) \, du \, ds = F_{(1,1)}(t),
		\]
		so $(k\ast a_1)(t) = (k\ast F_{(1,1)})(t)$. 
		
		For the second term, let $v_a(s) = \int_0^1 a_\phi(s)\,d\phi$. Then, by Lemma~\ref{lemma:delta_integral}, $v_a'(s)=(\delta a)(s)$. Integrating by parts we get 
		\begin{equation} \label{eq.kdela}
			\int_0^t k(t-s)(\delta a)(s)\,ds = k(0)v_a(t)-k(t)v_a(0)+\int_0^t k'(t-s)v_a(s)\,ds.
		\end{equation}
		Next, we find an expression for $a_\phi(s)$: by the definition of $a_\phi$, the fact $a=f_1$, and the definition of $F$, we get 
		\[ 
		a_\phi(s) = \int_{(s-\phi)^+}^s a(u) \, du = \int_{(s-\phi)^+}^s \int_{(u-1)^+}^u f(y) \, dy \, du = F_{(\phi, 1)}(s). \]
		Therefore 
		\[
		v_a(s)=\int_0^1 F_{(\phi, 1)}(s)\,d\phi,
		\]
		so as $F_{(\phi,1)}(0)=0$, $v_a(0)=0$. Using this and $k(0)=0$ in \eqref{eq.kdela} we get 
		\[
		\int_0^t k(t-s)(\delta a)(s)\,ds = (k'\ast v_a)(t) = \left(k'\ast \int_0^1 F_{(\phi, 1)}(\cdot)\,d\phi\right)(t), \quad t\geq 0.	
		\]
		Combining this formula with the fact that $k\ast a_1 = k\ast F_{(1,1)}$, we have that 
		\[
		I_1(t)= (k\ast F_{(1,1)})(t)+\left(k'\ast \int_0^1 F_{(\phi, 1)}(\cdot)\,d\phi\right)(t), \quad t\geq 0.
		\]
		Finally, combining this formula with \eqref{eq.I2final}, we see that $I=k\ast f=I_1+I_2$ has the claimed representation. 
	\end{proof} 

\end{document}
